\section{Scope, results, and conventions}
\label{scope:sec}

\subsection{The research questions being continued}

The source corpus is the manuscript in
\texttt{Analysis/Polylogarithms/docs/manuscript} and the six research
archives in \texttt{docs/incoming}, frozen at the revision on the title
page \cite{Repo}. A read-only extraction verified the Git blob identity
of 704 relevant text, code, data, bibliography, and archive files.
This is a source-integrity statement, not a claim to have verified every
theorem in those files. The mathematical audit is targeted to the
questions and premises used below.

Three developments in the incoming material make a further exact
calculus possible. First, periodic Stieltjes functions have been given
consistent distributional and pointwise finite-part conventions.
Second, nonlinear coordinate changes have been reduced to finite
residue expressions. Third, the shifted height-one harmonic family
has an exact Gamma-ratio representation with a finite Stirling
specialization. The present work uses those developments as
established inputs and extends the questions they explicitly leave
open \cite{IncomingCoincident,IncomingTransport,IncomingExact}.

The main outcomes and their precise scope are as follows.

\begin{longtable}{@{}>{\raggedright\arraybackslash}p{0.25\textwidth}
 >{\raggedright\arraybackslash}p{0.48\textwidth}
 >{\raggedright\arraybackslash}p{0.20\textwidth}@{}}
\toprule
Incoming question or target & Result established here & Location\\
\midrule
\endhead
Coincident products beyond degree two &
A finite local completion for every number of factors; a fully
explicit cubic global kernel at all argument orders &
Theorem \ref{cub:general-completion};
Proposition \ref{cub:global}\\[3pt]
Cubic digamma finite part &
An exact reduction to $\omega'''(0)$, where
$\omega(t)=T(t,t,t)$; also an exact evaluation of $\omega''(0)$ &
Theorem \ref{cub:diagonal}\\[3pt]
Contact terms at a shift collision &
One explicit Gamma--harmonic generator for every delta coefficient;
no undetermined lower delta derivatives &
Theorem \ref{cont:main}\\[3pt]
Coordinate-invariant singular germs &
Necessary and sufficient coefficient conditions for every prescribed
order of tangency &
Theorems \ref{invg:classification},
\ref{invg:stieltjes-threshold}\\[3pt]
Normal Tornheim derivatives at integer inner orders &
Finite depth-two reductions at all integer inner orders, and
depth-one closure at nonpositive orders with arbitrary weights &
Theorems \ref{sp:integer}, \ref{sp:rationaltheorem}\\[3pt]
Spectral direction dependence &
Constant, linear, and quadratic terms on every ray with
$(a+c)(b+c)\ne0$ &
Theorem \ref{sp:raytheorem}\\[3pt]
Independent harmonic exponents &
Complete Laurent principal parts and finite parts for two moving
exponents at every trailing depth and common shift &
Theorem \ref{h2:thm:direction}\\
\bottomrule
\end{longtable}

These statements answer several archived research questions, but they
do not all have the same kind of closure. The cubic finite part is
reduced to a specified multizeta derivative, not evaluated in ordinary
zeta constants alone. The positive-inner-order colored formulas close
in depth-two polylogarithms. The weighted nonpositive-inner-order
theorem closes in ordinary polylogarithms and their order derivatives.
The two-exponent harmonic theorem eliminates its holomorphic
remainder through the finite part, but does not eliminate the
remainder from every higher regular coefficient.

\subsection{Relation to the literature}

The analytic ingredients are classical. Hurwitz Fourier expansions
and product integrals are treated by Espinosa and Moll
\cite{EM2002,EM2006}. Bailey and Borwein systematically study
Mordell--Tornheim order derivatives and their polylogarithmic integral
representations \cite{BB2014}. Colored Tornheim partial fractions already appear
in Zhao \cite{sp:Zhao}. Borwein and Dilcher give direction-dependent
Tornheim values and first derivatives on the equal-inner-order family
\cite{BorweinDilcher}; related Witten values are discussed by Romik
\cite{sp:Romik}. Logarithmic finite parts and their change of variables
have an established theory \cite{Latt2011}. Meromorphic continuation
and Laurent algorithms for multiple zeta functions at integer points
are supplied by Akiyama--Egami--Tanigawa and
Matsumoto--Onozuka--Wakabayashi \cite{h2:AET,h2:MOW}.

Accordingly, this article makes no claim to have discovered finite
parts, the Tornheim--Hurwitz relation, directional regularization, or
multiple-zeta continuation. The contributions are the stated explicit
completion and contact generators, the invariant-subspace
classification, the asymmetric quadratic formula, and the finite
two-exponent harmonic coefficient rule relative to the audited corpus.
The elementary weighted identities are included as proved, useful
consequences of a general rational coefficient method.

\subsection{Notation and finite-part prescriptions}

We use
\begin{equation}\label{scope:stieltjes}
 \zeta(1+u,a)=\frac1u+
       \sum_{n\ge0}\frac{(-1)^n}{n!}\gamma_n(a)u^n,
 \qquad \gamma_0(a)=-\psi(a),
 \qquad \gamma_n=\gamma_n(1).
\end{equation}
Here $\psi=\Gamma'/\Gamma$, and $\gamma=\gamma_0$ is Euler's
constant. A prime on $\psi$ is an argument derivative. On $\zeta(s,a)$,
a prime or $[k]$ superscript denotes a derivative in $s$ unless an
argument derivative is explicitly written. Thus
$\zeta^{[k]}(s,a)=\partial_s^k\zeta(s,a)$ and
$\Li_s^{[k]}(z)=\partial_s^k\Li_s(z)$.
All integer summation bases are raised to complex powers using their
real logarithms. A common shift $a$ is initially in $\Re a>0$, with the
analytic logarithm in that half-plane.

The generalized harmonic numbers are
\[
 H_N^{(j)}=\sum_{\nu=1}^N\nu^{-j},\qquad H_N=H_N^{(1)},\qquad H_0=0.
\]
The rising factorial is $(a)_n=a(a+1)\cdots(a+n-1)$; the falling
factorial is $a^{\underline n}=a(a-1)\cdots(a-n+1)$.
The second-kind Stirling number is
$\left\{\begin{smallmatrix}n\\k\end{smallmatrix}\right\}$;
$s(n,k)$ denotes the signed first-kind Stirling number.
Coefficient extraction $[u^j]F$ always refers to an ordinary power
series, with any exponential-generating factorials displayed
separately.

For a function with a finite logarithmic singular expansion at zero,
\begin{equation}\label{scope:FP}
 \FP\int_0^1 f(x)\,dx
   =\CT_{\varepsilon\downarrow0}\int_\varepsilon^1 f(x)\,dx.
\end{equation}
The constant term means the coefficient of
$\varepsilon^0(\log\varepsilon)^0$ after discarding divergent powers
and logarithms. It is a coordinate-dependent prescription. A spectral
Laurent constant $\CT_{u=0}F(u)$ is a different operation; identifying
the two requires a proof of their local counterterms. Such a proof is
given in Section~\ref{cub:sec}.

For periodic distributions, $\delta$ is the unit Dirac mass at the
origin of $\mathbb R/\mathbb Z$, $D$ is the distribution derivative,
and $\check F(x)=F(-x)$. We normalize convolution so that
\[
 (\check F*G)(a)=\int_0^1 F(x)G(x+a)\,dx
\]
whenever ordinary integration is valid. A finite part of a pointwise
product is not defined by multiplying distributions. Section
\ref{cont:sec} explicitly compares two independently defined
extensions of the same off-collision function.

\paragraph{Why the directions matter.}
There is no single holomorphic Taylor germ of the uncolored Tornheim
function at $(0,0,0)$. For example,
\[
 T(t,t,t)\big|_{t=0}^{\mathrm{cont}}=\frac13,\qquad
 T(0,0,t)\big|_{t=0}^{\mathrm{cont}}=\frac5{12}.
\]
A notation such as $T''(0)$ without a specified one-variable
restriction is therefore insufficient. The same distinction occurs
when one exponent of a harmonic sum is fixed before differentiating
another. Every derivative formula below preserves its stated
restriction.
